% =====================================================================
%  PLANTILLA DE ENTREGA — Entregable 02
%  Sistemas Multimedia · Curso 2026/27 · Prof. Pedro Pernías
%  Compilador: pdfLaTeX (Overleaf por defecto)
%
%  Esqueleto SIN enunciados: las instrucciones, los ejemplos y las
%  rúbricas de cada apartado están en el documento del entregable.
%  Aquí solo escribís vuestra respuesta.
%
%  Podéis ampliar o reducir los espacios y añadir filas a las tablas:
%  las medidas son orientativas del tamaño de respuesta esperado.
%  NO cambiéis la cabecera ni el orden de las partes y apartados.
% =====================================================================

\documentclass[11pt,a4paper]{article}

\usepackage[utf8]{inputenc}
\usepackage[T1]{fontenc}
\usepackage{lmodern}
\usepackage{microtype}
\usepackage[spanish]{babel}

\usepackage[left=2.5cm,right=2.5cm,top=2.8cm,bottom=2.5cm]{geometry}
\usepackage{array}
\usepackage{booktabs}
\usepackage{enumitem}
\usepackage{graphicx}
\usepackage{needspace}
\usepackage{fancyhdr}
\usepackage[hidelinks]{hyperref}

% ------------------------------------------------------------
% Cabecera y pie
% ------------------------------------------------------------
\pagestyle{fancy}
\fancyhf{}
\renewcommand{\headrulewidth}{0.4pt}
\fancyhead[L]{\small Entregable 02 \,\textperiodcentered\, Prof: Pedro Pernías}
\fancyfoot[C]{\thepage}

% ------------------------------------------------------------
% Estructura mínima
% ------------------------------------------------------------
\newcommand{\parte}[1]{\par\vspace{1.3em}\noindent{\Large\bfseries #1}\par\vspace{0.4em}}
\newcommand{\apartado}[1]{\par\vspace{0.9em}\needspace{6\baselineskip}%
  \noindent{\large\bfseries #1}\par\vspace{0.3em}}
% Mantiene un rótulo pegado a su tabla.
\newcommand{\bloque}[2]{\par\noindent\begin{minipage}{\linewidth}%
  \noindent#1\par\vspace{0.3em}#2\end{minipage}\par}

% ------------------------------------------------------------
% MACRO FIJA: tabla de integrantes (hasta 6). No la modifiquéis.
% ------------------------------------------------------------
\newcommand{\integrantes}{%
  \noindent\begin{tabular}{c p{7.6cm} p{5.6cm}}
  \toprule
  \textbf{\#} & \textbf{Apellidos} & \textbf{Nombre} \\
  \midrule
  1 & \rule{0pt}{2.0em} & \\
  \midrule
  2 & \rule{0pt}{2.0em} & \\
  \midrule
  3 & \rule{0pt}{2.0em} & \\
  \midrule
  4 & \rule{0pt}{2.0em} & \\
  \midrule
  5 & \rule{0pt}{2.0em} & \\
  \midrule
  6 & \rule{0pt}{2.0em} & \\
  \bottomrule
  \end{tabular}}

\setlist{itemsep=2pt, topsep=3pt}

\begin{document}

% =====================================================================
%  CABECERA DE ENTREGA
% =====================================================================
\noindent{\LARGE\bfseries Entregable 2: El Lean Startup Canvas reducido de vuestro producto}\par
\vspace{0.6em}
\noindent\textbf{Grupo:} \hrulefill \quad \textbf{Producto o \emph{startup}:} \hrulefill

\vspace{0.8em}
\noindent\textbf{Integrantes del grupo (hasta 6):}\par\vspace{0.4em}
\integrantes

\parte{PARTE A. Cliente y problema}

\apartado{A.1. Segmento de clientes}
\vspace{4.0cm}

\apartado{A.2. Problema}
\noindent\begin{tabular}{p{0.8cm} p{6.6cm} p{6.4cm}}
\toprule
\textbf{\#} & \textbf{Problema (por orden de importancia)} & \textbf{Alternativa que usa hoy} \\
\midrule
1 \rule{0pt}{2.6em} &  &  \\
\midrule
2 \rule{0pt}{2.6em} &  &  \\
\midrule
3 \rule{0pt}{2.6em} &  &  \\
\bottomrule
\end{tabular}

\parte{PARTE B. Solución, propuesta de valor y ventaja}

\apartado{B.1. Solución}
\noindent\begin{tabular}{p{0.8cm} p{6.6cm} p{6.4cm}}
\toprule
\textbf{\#} & \textbf{Característica} & \textbf{Problema al que responde} \\
\midrule
1 \rule{0pt}{2.6em} &  &  \\
\midrule
2 \rule{0pt}{2.6em} &  &  \\
\midrule
3 \rule{0pt}{2.6em} &  &  \\
\bottomrule
\end{tabular}

\apartado{B.2. Propuesta de valor única}
\vspace{3.0cm}

\apartado{B.3. Ventaja especial (o ``injusta'')}
\vspace{4.0cm}

\parte{PARTE C. Riesgo clave y validación inicial}

\apartado{C.1. Riesgo clave}
\vspace{5.0cm}

% =====================================================================
%  DECLARACIÓN DE USO DE LA IA
% =====================================================================
\parte{Declaración de uso de la IA}
\vspace{5.0cm}

\end{document}
